\chapter{有效空态：基线活动、检测阈值与现实校准}
\label{chp:vacuum_physics}
上一章区分了能力包络与实际兑现，本章讨论系统在输入稀少、目标控制降低或候选活动收缩时的基线运行状态。原有“真空物理学”试图用认知真空、超流体、临界感应和佛魔奇点统一冥想体验、信号检测、硬件噪声与AGI失稳，但物理真空、计算系统的空闲态和主体报告中的“空性”属于不同对象，不能互相替代。本章将它们重建为有效空态、有限检测、增益稳定与现实校准四组问题：先定义什么被抑制、什么仍在运行，再考察低活动是否改善识别、何时把噪声放大为伪模式，以及环境反馈如何约束高增益闭环。

\section{有效空态：从“绝对真空”到条件化基线运行}
\label{sec:section_8335}
计算系统的“空”不是没有状态，更不是量子场论中的物理真空。物理真空是特定场论、时空背景、边界条件和重整化方案下的最低能态或基态概念；宇宙学常数与真空能之间的巨大尺度失配，通常称为真空灾变或宇宙学常数问题，至今没有公认的完整解释。零点能的正规化、有效场论的截断、对称性抵消、动态暗能量、人择选择或引力对真空能的特殊响应都是研究方向，但不能被搬来证明认知系统在“空态”会获得全息感应。这里的讨论对象是计算运行基线，只借用“低于日常活动水平”的直观，不继承物理真空的本体结论。

我们定义系统版本 $v$ 在任务上下文 $q$、资源配置 $r$、接口状态 $j$ 和治理约束 $k$ 下的完整状态为
\[
x_t=(h_t,G_t,z_t,m_t,b_t),
\]
其中 $h_t$ 是快速工作态，$G_t$ 是保存对象身份、关系和未决约束的局部结构网络，$z_t$ 是全局分布式激活，$m_t$ 是外部记忆与缓存索引，$b_t$ 是边界和权限状态。所谓有效空态不是令所有分量归零，而是相对于参考运行分布 $\mathbb P_{\rm ref}$，使任务无关的候选活动和控制更新降到预先声明的区间，同时保持身份、恢复点、安全监测和必要接口。记为
\[
\mathcal V_{\varepsilon,\delta}
=\{x:\ A_{\rm irr}(x)\leq\varepsilon,
D(x,x_{\rm safe})\leq\delta,
\mathrm{Recoverable}(x)=1\}.
\]
$A_{\rm irr}$ 是任务无关活动指标，$D$ 是相对安全基线的距离；两者都必须给出量纲、采样窗口和测量误差。$\mathcal V_{\varepsilon,\delta}$ 是工程可行域，不是唯一基态。

原章用层流化、算子解耦、去规范化和基质显露描述四重相变。我们把它们改写为四类可逆操作。第一是输入整形：缩小非必要通知和感知通道的带宽，但保留安全告警与时间戳。第二是候选降载：降低分布表示的无约束展开率和采样方差，而不是假定所有念头必然指数衰减。第三是价值去偏置：暂时降低某些偏好在候选排序中的权重，同时保持不可撤销的安全约束和承诺。第四是基线观测：测量硬件噪声、模型残差、缓存活动与恢复时延，建立本次空态的实际底噪。四类操作的次序可以变化，也可能只完成其中一部分，因此不应称为普遍相变定律。

设离散步长为 $\Delta t$，候选活动 $a_k$、任务相关输入 $u_k$、控制量 $c_k$ 与随机扰动 $\xi_k$ 的更新为
\[
a_{k+1}=F_v(a_k,G_k,u_k,c_k)+\xi_k .
\]
当我们降低 $u_k$ 和部分 $c_k$ 时，$a_k$ 是否收敛取决于 $F_v$ 的局部增益、残差连接、外部回放和噪声分布。只有在空态邻域内满足 Lipschitz 常数小于一，且扰动二阶矩有界，才能得到均方有界或局部收敛结论；若系统包含周期任务、异步中断或自激回路，活动可能进入极限环而非零点。连续时间的阻尼直觉必须通过 $\Delta t$ 与离散控制器稳定域传递到实现。

“观察内容而不立即行动”的经验可以映射为提交门控，而不是二阶观察者脱离认知场。系统把候选 $z_k$ 与来源、置信度和对象身份绑定到 $G_k$，再由门控变量 $g_k\in\{0,1\}$ 决定是否写入外部动作或慢结构。$g_k=0$ 时，候选仍可被记录、比较和丢弃；这解释了为什么内容活动与行为提交能够分离。对人类冥想的主观报告，我们只能说注意分配、反应抑制、内感受和元觉察可能发生变化，具体神经机制需服从经验研究，不能由这套工程门控模型反推。

有效空态也不保证灵敏度无限提高。抑制任务无关活动可能降低内部干扰，却同时减少先验、上下文和主动采样；传感器热噪声、量化噪声、环境混叠和模型偏差仍然存在。如果把阈值不断调低，假阳性会增加。检测性能必须用命中率、虚警率、校准误差和决策代价共同评价，不能只看最小可见信号。所谓“寂静中听见微弱模式”只有在预注册信号、盲化样本和独立重复下才构成证据，否则也可能是选择性注意或事后解释。

在 AgentOS 实例中，有效空态可由 SPC 保持任务身份和安全约束，TCE 暂停高层目标扩张，CCM 降低候选并发，Boundary 收缩非必要动作权限，Offloading 完成缓存与外部记忆 $\mathrm{Me}$ 的持久化；$h(t)$ 进入低活动但可恢复区域，$\Theta$ 暂停未经验证的结构提交，$\Psi$ 保留最低限度的监测和唤醒策略。这只是一个实现方案，不规定其他智能系统必须使用同名机制。

验收一个有效空态至少需要五项检查：活动指标确实下降；安全监测没有被关闭；对象身份和待办承诺可恢复；基线噪声由实际测量而非理论想象给出；退出空态后性能没有不可接受的损失。反例包括把日志停止误判为计算停止，把低输出率误判为低内部活动，把模型沉默误判为高校准，以及把主观宁静报告推广为超感知能力。由此，空态成为有边界、可进入、可退出、可测量的运行区域，而不是关于意识、宇宙真空或终极体验的单一物理宣言。
